\documentclass[12pt]{article}
\usepackage[dvips]{graphicx}
\usepackage{setspace,palatino,multirow}
\usepackage{amsmath,amssymb}
\usepackage{amsthm}
\usepackage{subfig,comment}
\usepackage{pdflscape}
\usepackage[dvipsnames]{xcolor}
\usepackage{hyperref}
\usepackage[longnamesfirst]{natbib}
\bibliographystyle{chicago}

\special{papersize=8.5in,11in}

\definecolor{darkred}{rgb}{0.66,0,0}
\definecolor{darkblue}{rgb}{0,0,0.25}

\hypersetup{
  colorlinks=true,
  citecolor=darkblue,
  linkcolor=darkblue,
  pdfpagemode=none,
  pdfstartview={FitH},
  pdftitle={},
  pdfauthor={},
  pdfkeywords={}
}

\setlength{\parindent}{0.25in}
\setlength{\parskip}{1.5ex plus0.5ex minus0.5ex}
\textwidth15.75cm
\evensidemargin5mm
\oddsidemargin5mm
\topmargin-15mm
\textheight 22.7cm
\hyphenation{over-lapping}

\renewcommand{\arraystretch}{1.1}
\usepackage[margin=2cm]{geometry}

\newcounter{saveeqni}%
\newcounter{saveeqn01i}%
\newcommand{\alpheqni}{\setcounter{saveeqni}{\value{section}}%
\renewcommand{\theequation}
    {\alph{saveeqni}\mbox{.\arabic{equation}}}}%
\newcommand{\reseteqni}{\setcounter{equation}{\value{saveeqni}}%
\renewcommand{\theequation}{\arabic{equation}}}%

\begin{document}
\theoremstyle{definition}
\newtheorem{definition}{Definition}
\theoremstyle{definition}
\newtheorem{result}{Result}
\theoremstyle{plain}
\newtheorem{lemma}{Lemma}
\theoremstyle{plain}
\newtheorem{proposition}{Proposition}
\theoremstyle{plain}
\newtheorem{theorem}{Theorem}
\theoremstyle{plain}
\newtheorem{corollary}{Corollary}

\begin{onehalfspacing}

\vspace{1.0cm}

\noindent \textbf{Reply to Referee 1 of the manuscript ``Resolving Coordination Frictions in Green Labor Transitions: Minimizing Unemployment, Costs, and Welfare Distortions.''}

\bigskip

Dear Referee,

\medskip

Thank you very much for your careful reading and highly constructive report. I am grateful for both your two essential comments and your detailed smaller suggestions. I revised the manuscript to address all points in a systematic way.

In what follows, I reproduce each comment in bold and provide my response directly below it.

Before turning to the comment-by-comment responses, I summarize the main revisions:

\textbf{First}, I strengthened the transition section by clarifying what is and is not claimed under the sequence-of-steady-states approach, and by adding quantitative diagnostics to evaluate transition speed after policy shocks.

\textbf{Second}, I expanded the calibration/identification discussion with a clearer parameter-to-moment mapping and additional intuition for why key moments are most informative for specific parameters.

\textbf{Third}, I implemented all editorial and presentation changes requested (terminology, figure presentation, framing, and section-level clarity), and clarified several ambiguous statements in the calibration and policy discussion.

I now move to the specific comments.

\bigskip
\bigskip

\textbf{Comment 1 (Essential): Transition dynamics. I am satisfied with the way the paper models the green transition as a sequence of steady states. That said, this will seem more plausible if the paper can show that the economy transitions quickly to a new steady state after unexpected policy shifts (e.g., MIT shocks).}

\medskip

Thank you for this important suggestion. I agree this is central for interpretation. I now add a quantitative transition-speed exercise after unanticipated policy changes and report convergence metrics (including [INSERT METRICS]). These results support the practical use of a sequence-of-steady-states approximation in my setting.

\bigskip
\bigskip

\textbf{Comment 2 (Essential): Identification. The paper maps parameters to targeted moments, but needs a clear and intuitive explanation of this mapping. It would help to show how key moments vary when key parameters change.}

\medskip

Thank you for this valuable point. I expanded the identification discussion and now provide an intuitive mapping from parameters to moments, including why some moments are more sensitive to specific parameters. I also add sensitivity evidence (tables/figures) showing how key targeted moments move with parameter variation.

\bigskip
\bigskip

\textbf{Comment 3 (Smaller): Clarify that a hiring likelihood ratio of 1.29 means ``29\% higher than for other workers.''}

\medskip

Implemented. I now state this explicitly in the calibration paragraph and align wording with the value in the table.

\bigskip
\bigskip

\textbf{Comment 4 (Smaller): The statements ``using 2022 generation is conservative given the near-term pipeline'' and ``the higher benchmark provides a fiscal stress test'' are unclear.}

\medskip

Thank you. I rewrote this passage to define both phrases precisely and to explain the logic in transparent, quantitative terms.

\bigskip
\bigskip

\textbf{Comment 5 (Smaller): Why does calibration use a 2023 unemployment rate of 3.5\% while claiming calibration to the 2022 labor market?}

\medskip

Thank you for catching this inconsistency. I corrected and clarified the date convention in the calibration section and ensured internal consistency across text, targets, and tables.

\bigskip
\bigskip

\textbf{Comment 6 (Smaller): Under Table 3, ``three policy paths'' should be ``three policy mixes'' (or equivalent), since these are not paths yet.}

\medskip

Implemented. I replaced this language and now consistently refer to these as policy mixes/configurations at a fixed target.

\bigskip
\bigskip

\textbf{Comment 7 (Smaller): In Figure 7, why are bars not the same width?}

\medskip

Thank you. I revised the figure rendering to ensure consistent bar widths and improved visual readability.

\bigskip
\bigskip

\textbf{Comment 8 (Smaller): The key takeaway from Figure 7 appears to be the strong rotation in policy mix versus relatively stable total intervention size; hump-shape seems second-order.}

\medskip

I agree and revised the text accordingly. The discussion now emphasizes composition rotation as the primary mechanism/result, with the hump profile presented as a secondary quantitative feature.

\bigskip
\bigskip

\textbf{Comment 9 (Smaller): Writing in Section 4.3.1 is somewhat convoluted.}

\medskip

Thank you. I rewrote this subsection for clarity, shortening dense passages and improving signposting of the economic mechanism.

\bigskip
\bigskip

\textbf{Comment 10 (Smaller): Policy intervention sizes are likely an upper bound because technical change may raise green productivity/wages and younger cohorts may face lower entry costs over time.}

\medskip

I agree and now add explicit upper-bound framing and interpretation in the discussion/conclusion, including why this does not overturn the core instrument-design ranking.

\bigskip
\bigskip

\textbf{Comment 11 (Smaller): Figure D.8 in the appendix is very illustrative and may deserve to be in the main text.}

\medskip

Thank you for this suggestion. I moved [INSERT FIGURE ID] to the main text (or cross-referenced it more prominently, depending on space constraints) and updated the surrounding discussion.

\bigskip

I again thank you for the insightful and constructive feedback. It substantially improved both the clarity and the policy relevance of the revised manuscript.

\end{onehalfspacing}
\end{document}
